\documentclass[11pt]{article}
\usepackage[a4paper,margin=27mm]{geometry}
\usepackage[T1]{fontenc}
\usepackage{lmodern,amsmath,amssymb,amsthm,mathtools,microtype,booktabs,tabularx}
\usepackage[colorlinks=true,linkcolor=blue,urlcolor=blue]{hyperref}
\newcommand{\D}{\mathbb D}
\newcommand{\E}{\mathbb E}
\newcommand{\Pp}{\mathbb P}
\newcommand{\dd}{\,d}
\newtheorem{conjecture}{Conjecture}
\title{GMC representations and the intensity conjecture\\[3pt]\large A side note on construction versus selection}
\author{}
\date{8 September 2026}
\begin{document}
\maketitle
\vspace{-1.8em}
\noindent\textbf{Position.}
The Gaussian multiplicative chaos (GMC) representation does not, by itself,
change the leading conjecture
\begin{equation}\label{target}
 \rho_\beta(z)=\frac{4/\beta-1}{\pi(1-|z|^2)^2}.
\end{equation}
It provides an explicit family of hyperbolically invariant processes,
including a candidate with this intensity, and exposes a precise selection
problem. Constructing that candidate is not a proof that the finite disk gas
converges to it. This note records those distinctions. It uses the random
Bergman construction developed in the preceding discussion as an input;
it does not reproduce its analytical proof.

\section{The finite gas uses a biased Gaussian law}
Let $\mu=dA/\pi$ on the unit disk, and consider the labelled finite gas
\[
 d\Pp_{N,\beta}(\mathbf z)
 =Z_N(\beta)^{-1}|\Delta_N(\mathbf z)|^\beta\,d\mu^N(\mathbf z),
 \qquad \Delta_N=\prod_{i<j}(z_i-z_j).
\]
Throughout this note, $0<\beta<2$ and $t=2-\beta\in(0,2)$.
Take a centered real Gaussian generalized field $X$ with covariance
$C(z,w)=-\log|z-w|$, and let
\[
 M_t^X(dz)={:}e^{\sqrt t X(z)}{:}\,d\mu(z).
\]
The colons denote Wick normalization in the GMC limit. For a finite measure
$M$ with full support, define
\[
 G_N(M)_{jk}=\int z^j\overline z^{\,k}\,M(dz),\quad 0\le j,k<N,
 \qquad D_N(M)=N!\det G_N(M).
\]
The determinant identity and the GMC product-moment identity give
\[
 D_N(M)=\int|\Delta_N|^2\,dM^{\otimes N},\qquad
 \E D_N(M_t^X)=Z_N(2-t).
\]
Hence the exact mixing law for the original gas is
\begin{equation}\label{bias}
 \frac{dQ_{N,t}}{dP_{\rm GMC}}(M)
 =\frac{D_N(M)}{\E D_N(M_t^X)}.
\end{equation}
Conditionally on $M$, sample the rank-$N$ polynomial projection DPP in
$L^2(M)$. Its labelled law is
$|\Delta_N|^2\,dM^{\otimes N}/D_N(M)$. Averaging with \eqref{bias}
cancels $D_N$ and recovers $\Pp_{N,2-t}$ exactly.

In contrast, averaging the same conditional DPP with the original Gaussian
law $P_{\rm GMC}$ defines the \emph{unbiased mixture}. Its asymptotics do
not automatically describe the gas. Although $D_N$ is a normalization
factor inside the conditional particle law, it is a nonconstant random
weight in the law of the environment.

\section{What the infinite GMC construction supplies}
Let $G$ be the zero-boundary Gaussian free field with covariance
\[
 g_\D(z,w)=-\log\left|\frac{z-w}{1-z\overline w}\right|,
 \qquad M_t^G={:}e^{\sqrt tG}{:}\,d\mu.
\]
For $s>-1$, put
\[
 m_{t,s}(dz)=(1-|z|^2)^sM_t^G(dz).
\]
Write $\mathcal Y_{t,s}$ for the following position law: first sample
$M_t^G$, then sample the infinite projection DPP corresponding to the
closure of $\mathbb C[z]$ in $L^2(m_{t,s})$.

\noindent\textbf{Construction result used here.}
For each fixed $t\in(0,2)$ and $s>-1$, this projection is locally trace
class almost surely; its finite polynomial projections converge locally.
The annealed process $\mathcal Y_{t,s}$ is locally finite, almost surely
infinite, and hyperbolically invariant, with ordinary-area intensity
\begin{equation}\label{family}
 \rho_{t,s}^{\mathcal Y}(z)
 =\frac{s+1}{\pi(1-|z|^2)^2}.
\end{equation}
The coefficient is independent of $t$ at fixed $s$; the entire law is not
asserted to be independent of $t$. The proof uses bounded analytic point
evaluations from GMC small-ball estimates, followed by an equivariant
polar decomposition and hyperbolic mass transport. Standard negative
GMC moments are one input; see \cite{negative}.

To relate this family to the unbiased mixture, decompose
\[
 X=G+\operatorname{Re}F,
 \qquad \E F(z)\overline{F(w)}=-2\log(1-z\overline w),
\]
where $G$ and the analytic Gaussian field $F$ are independent. Wick
normalization gives
\[
 M_t^X=|h|^2(1-|z|^2)^{t/2}M_t^G,
 \qquad h=e^{\sqrt t F/2}.
\]
The harmonic factor can be removed from the infinite polynomial DPP.
Precisely, with $m_0=(1-|z|^2)^{t/2}M_t^G$,
\[
 \overline{h\mathbb C[z]}^{\,L^2(m_0)}
 =\overline{\mathbb C[z]}^{\,L^2(m_0)}\quad\text{almost surely}.
\]
This density statement is essential because $h$ need not be bounded.
It follows by approximating $h$ and $1$ using $h(rz)$ and $h(z)/h(rz)$
in $L^2(m_0)$ along a suitable sequence $r\uparrow1$.
Multiplication by $h$ then gives a unitary identification, and its kernel
factors cancel the change of reference measure in every DPP correlation
measure. Thus the unbiased infinite mixture is $\mathcal Y_{t,t/2}$, with
\begin{equation}\label{unbiased}
 \rho_\beta^{\rm unbiased}(z)
 =\frac{1+t/2}{\pi(1-|z|^2)^2}
 =\frac{2-\beta/2}{\pi(1-|z|^2)^2}.
\end{equation}

\section{Construction does not verify the conjecture}
Let $b$ denote the coefficient multiplying
$1/[\pi(1-|z|^2)^2]$. The relevant comparison is
\begin{center}
\begin{tabularx}{\linewidth}{@{}Xcc@{}}
\toprule
Position law & Exponent $s$ & Coefficient $b$\\
\midrule
Unbiased GMC limit & $t/2$ & $1+t/2$\\[3pt]
GMC candidate matching \eqref{target} & $2t/(2-t)$ & $(2+t)/(2-t)$\\[3pt]
Actual disk-gas limit & Unidentified & Unproved at fixed $\beta\ne2$\\
\bottomrule
\end{tabularx}
\end{center}
The actual limit is not assumed to belong to the $\mathcal Y_{t,s}$ family.
The two explicit coefficients satisfy
\[
 \frac{b_{\rm conjectured}}{b_{\rm unbiased}}=\frac2\beta,
 \qquad
 b_{\rm conjectured}-b_{\rm unbiased}
 =\frac{(2-\beta)(4-\beta)}{2\beta}>0.
\]
Therefore removing the determinant weight cannot establish the proposed
coefficient. This comparison does \emph{not} prove that the actual
fixed-temperature gas limit differs from the unbiased mixture: its
intensity is precisely what remains unknown.

More importantly, \eqref{family} realizes \emph{every} positive
coefficient by choosing $s=b-1$. The choice
\begin{equation}\label{sstar}
 s_*(\beta)=\frac4\beta-2=\frac{2(2-\beta)}\beta
\end{equation}
is obtained by solving $s+1=4/\beta-1$. It builds the proposed intensity
into a candidate; it is not an independent derivation of that intensity
from the finite gas. Existence of many invariant candidates also does not
show nonuniqueness of the finite-$N$ limit, or nonuniqueness for one fixed
DLR specification. A common specification has not been established for
this family.

\section{Perturbative evidence and its limitation}
In terms of $t=2-\beta$, the leading conjecture predicts
\[
 b_{\rm conjectured}=\frac{2+t}{2-t}
 =1+t+\frac{t^2}{2}+O(t^3),
 \qquad
 b_{\rm unbiased}=1+\frac t2.
\]
The previously obtained first temperature response of the finite gas is
\[
 \lim_{N\to\infty}
 \left.\partial_\beta\rho_{N,\beta}(z)\right|_{\beta=2}
 =-\frac1{\pi(1-|z|^2)^2},
\]
locally uniformly in the disk. The earlier second-response calculation
likewise matches the quadratic term of the first expansion. These are
inputs from the preceding calculations, not new derivative results of
this side note.

They support retaining \eqref{target} as the leading conjecture. However,
derivatives were taken \emph{before} $N\to\infty$. Without a justified
exchange of limits or a uniform temperature remainder, those results do
not exclude the unbiased coefficient as a possible fixed-temperature
limit. Pointwise convergence of analytic functions on a real interval
alone does not imply convergence of their derivatives.

\section{A sharper conjecture to test}
\begin{conjecture}[Tentative GMC selection conjecture]\label{selection}
For each $0<\beta<2$, the point process of the original $N$-particle disk
gas converges locally in distribution to
\[
 \mathcal Y_{\,2-\beta,\,2(2-\beta)/\beta}.
\]
\end{conjecture}
This is stronger than the intensity conjecture: it specifies the entire
limiting law. If true, it would give full-sequence convergence, hyperbolic
invariance, and \eqref{target} for the limiting process. Its plausibility
must be tested independently through higher correlations, Palm laws, and
conditional resampling laws, including the probabilities of different
particle counts. Matching an intensity is insufficient. In particular,
we have not proved that this candidate satisfies the required limiting
DLR specification.

One exact way to formulate the remaining selection mechanism is through
orthogonal-polynomial norms. Let
\[
 h_N(M)=\min_{\substack{p\text{ monic}\\\deg p=N}}\int|p|^2\,dM.
\]
Since $D_N=N!\prod_{j=0}^{N-1}h_j$, the biased environment laws obey
\begin{equation}\label{increment}
 \frac{dQ_{N+1,t}}{dQ_{N,t}}(M)
 =\frac{h_N(M)}{\E_{Q_{N,t}}h_N(M)}.
\end{equation}
Identifying the effect of these successive weights on local particle
observables is a concrete unresolved task. Fixed-environment convergence
of the conditional DPP does not justify replacing $Q_{N,t}$ by the
unbiased law.

The original leading conjecture is therefore retained, while
Conjecture~\ref{selection} is an additional, more tentative proposal.
All constructions in this note concern $0<\beta<2$; the positive real
GMC representation used here gives no direct conclusion for $\beta>2$.

\begin{thebibliography}{9}
\bibitem{negative}
C. Garban, N. Holden, A. Sep\'ulveda and X. Sun,
\emph{Negative moments for Gaussian multiplicative chaos on fractal sets},
Electron. Commun. Probab. 23 (2018), no.~100.
\href{https://arxiv.org/abs/1805.00864}{arXiv:1805.00864}.
\end{thebibliography}
\end{document}
